% FILE: 06_contribution_to_thesis.tex
% Project: schemas
% Thesis Role: data-schema-definitions-for-ma-and-receipt-surfaces

\section{Contribution to the Dissertation}
\label{sec:schemas-contribution}

\subsection{Contribution to Prediction vs.~Coordination}
\label{sec:schemas-contribution-prediction}

The central distinction drawn in this dissertation is between systems that predict process behavior
and systems that coordinate process evidence. Prediction systems consume event logs and produce
probability distributions, forecasts, or recommendations. Coordination systems consume evidence
artifacts and produce admissible claims---claims that a counterparty can independently verify.

The schemas project contributes to the coordination side of this distinction by defining the
contract that makes coordination possible. A prediction system can produce a fitness score of 0.982
without any structural requirement that the score be linked to a specific event log by content
address. A coordination system, operating under the \texttt{SlideVerificationLedger} schema,
must produce a receipt that carries a 64-hex BLAKE3 hash of the event log, a 64-hex BLAKE3 hash
of the process model, a query definition identifying the execution engine, and an Ed25519 signature
over the JCS-canonical payload.

This difference is not cosmetic. It determines whether a claim can be independently replayed. A
prediction system's output is opaque to a counterparty who does not have access to the original
model and data. A coordination system's receipt is transparent: it names the inputs by content
address, names the execution engine, and certifies the result with a verifiable signature. The
schema is what enforces this transparency as a structural property rather than a voluntary practice.

The dissertation contribution is therefore: a receipts-based coordination architecture requires
a schema layer that elevates receipt transparency from convention to contract.

\subsection{Contribution to Process Evidence}
\label{sec:schemas-contribution-evidence}

Process evidence, as defined in this research program, is evidence that a lawful process happened.
It is not evidence of a model that could describe a process, or a metric that summarizes a process,
or a dashboard that visualizes a process. It is a traceable chain from a specific event log,
through a specific model, via a specific query, to a specific result---signed by a verifiable key.

The schemas project contributes to the process evidence architecture by specifying the shape of
the terminal artifact in this chain: the receipt. Without the schema, the chain terminates in an
informal document. With the schema, the chain terminates in a typed artifact that a validator can
accept or reject without consulting the original analyst.

The \texttt{OTelTraceBlake3IntegrationSchema} extends this contribution to the execution layer.
By requiring a per-span \texttt{blake3\_receipt} computed as
$\operatorname{BLAKE3}(\text{PriorReceipt} \| \text{SpanData})$, the schema makes the execution
path itself part of the evidence chain. An auditor can verify not only that the result is as
claimed, but that the execution that produced it proceeded through a specific sequence of spans
in a specific order with specific instruction counts. This is a stronger form of process evidence
than conformance checking alone provides.

The contribution to the dissertation is: process evidence requires two schema layers---one
governing the content of the analysis inputs (log hash, model hash), and one governing the
structure of the execution trace (span chain, BLAKE3 receipts per span).

\subsection{Contribution to Manufacturing Architecture}
\label{sec:schemas-contribution-manufacturing}

The manufacturing architecture of this research program---as expressed in the Chatman Equation
$A = \mu(O^*, T, L)$---requires a type law surface $T$ that constrains what artifacts are
admissible. The schemas project materializes the type law surface for the receipt domain.

This contribution is architectural in a specific sense: the schema is not a test, not a policy,
and not a recommendation. It is a structural constraint that the manufacturing function $\mu$ must
satisfy. Any receipt that does not conform to the \texttt{SlideVerificationLedger} schema is not
a valid output of the manufacturing function---it is a malformed artifact that the proof gate must
refuse.

The \texttt{additionalProperties: false} constraint is the architectural expression of this
principle. It means the schema does not merely describe what a receipt should contain; it refuses
what a receipt must not contain. This is the difference between a guideline and a law.

Within the seven-step M\&A deck rendering pipeline, the schema governs Step~3 (Cross-Reference
Compat Evidence Ledger) and Step~6 (Validate Artifacts). Its presence in both steps reflects its
role as both an input contract (Step~3: receipts must conform before they are cross-referenced)
and an output contract (Step~6: receipts must be verified against the schema before the deck is
sealed).

\subsection{Contribution to Capital-Flow Infrastructure}
\label{sec:schemas-contribution-capital}

The capital-flow and settlement use case for process intelligence---including M\&A transactions,
DAO governance, and board-level financial claims---requires that process evidence carry legal
defensibility. Legal defensibility requires that claims be independently verifiable, which requires
that the artifacts underlying the claims be identifiable by content address and signed by a
verifiable key.

The schemas project is the lowest-level technical artifact that enables this capital-flow
infrastructure. The five production receipts demonstrate the application to concrete M\&A claims:
\$1,250,000 EBITDA improvement, \$1,369,863 working capital release, \$450,000 SLA cap,
\$0 compliance leakage, and 97.5\% conformance standardization. Each of these dollar-value
assertions is backed by a receipt that identifies its evidential basis by BLAKE3 hash, names its
execution engine, and carries an Ed25519 signature.

The slide-to-receipt map (\texttt{ma/slide-to-receipt-map.md}) operationalizes this capital-flow
connection by linking each M\&A deck slide to its receipt file with explicit dollar-value
assertions and defensibility rules. The schema is what makes this linkage mechanically verifiable
rather than informally asserted. Without the schema, the slide-to-receipt map is a documentation
convention. With the schema, it is an enforceable contract.

The contribution to capital-flow infrastructure is: a receipts schema is the foundational artifact
that makes process intelligence claims legally defensible at the structural level---prior to any
question of whether the underlying process mining analysis was conducted correctly.
